\documentclass{article}
\usepackage{enumitem}
\usepackage{amsmath}
\begin{document}

\title{Chapter 37: Biodiversity and Conservation}
\date{}
\maketitle

\begin{enumerate}[label=Q\arabic*.]

\item According to the IUCN Red List, a species is classified as ``critically endangered'' when:
\\(A) Population has declined by 20\% in 10 years
\\(B) It faces an extremely high risk of extinction; population may have declined by $\geq$80\% over 10 years
\\(C) It is found in only one country
\\(D) Its population is below 10,000 individuals

\item The HIPPO acronym summarizes the major causes of biodiversity loss. What does it stand for?
\\(A) Habitat destruction, Invasive species, Pollution, Population (human), Over-exploitation
\\(B) Heat, Irrigation, Pesticide, Population, Overhunting
\\(C) Habitat loss, Inbreeding, Pathogens, Predation, Over-harvesting
\\(D) Hybridization, Invasives, Pollution, Poaching, Overgrazing

\item Biodiversity hotspots are regions that:
\\(A) Have the highest number of endemic species AND have lost $>$70\% of their original habitat
\\(B) Have the highest total species richness
\\(C) Are protected by international law
\\(D) Have the highest productivity

\item Which of the following is an example of in situ conservation?
\\(A) Cryopreservation of seeds in gene banks
\\(B) Zoological parks and botanical gardens
\\(C) National parks and biosphere reserves
\\(D) Tissue culture of endangered plants

\item Assertion: Biodiversity loss reduces ecosystem resilience.\\
Reason: Greater species diversity provides functional redundancy, buffering ecosystems against disturbances and species loss.
\\(A) Both true, Reason is correct explanation
\\(B) Both true, Reason is not correct explanation
\\(C) Assertion true, Reason false
\\(D) Assertion false, Reason true

\item The Convention on Biological Diversity (CBD, 1992, Rio de Janeiro) recognized:
\\(A) That developed nations own global biodiversity
\\(B) Sovereign rights of nations over biological resources and need for equitable sharing of benefits from biodiversity utilization
\\(C) That all species should be brought into captivity for protection
\\(D) Climate change as the primary threat to biodiversity

\item Species diversity includes:
\\(A) Only the number of species (species richness)
\\(B) Both species richness and relative abundance (evenness) of species
\\(C) Only genetic diversity within a species
\\(D) Diversity of ecosystems in a region

\item The ``rivet popper'' hypothesis (Ehrlich) for the importance of biodiversity states:
\\(A) Species can be removed from an ecosystem without harm until a threshold
\\(B) Like rivets in an airplane, each species' loss weakens the ecosystem; loss of any species increases probability of ecosystem collapse
\\(C) Only keystone species matter for ecosystem function
\\(D) Biodiversity is more important in tropical than temperate ecosystems

\item Biosphere reserves differ from national parks in that:
\\(A) Biosphere reserves exclude human populations entirely
\\(B) Biosphere reserves have a zonation system (core, buffer, transition zones) allowing human use in outer zones while strictly protecting the core
\\(C) National parks include marine areas; biosphere reserves are only terrestrial
\\(D) Biosphere reserves focus only on a single endangered species

\item Ex situ conservation strategies include all EXCEPT:
\\(A) Seed banks
\\(B) Zoological parks
\\(C) Wildlife sanctuaries
\\(D) Cryopreservation of gametes

\item The latitudinal diversity gradient (higher biodiversity near the equator) is explained by:
\\(A) More precipitation in tropical regions only
\\(B) Greater evolutionary time, higher productivity, more stable climate, and greater niche complexity in tropics
\\(C) Fewer predators in tropical regions
\\(D) Less competition in tropical ecosystems

\item Keystone species are defined as:
\\(A) The most abundant species in an ecosystem
\\(B) Species with disproportionately large effects on ecosystem structure relative to their abundance
\\(C) Species at the top of the food chain
\\(D) Endemic species unique to a region

\item Project Tiger was launched in India in 1973 to:
\\(A) Promote tiger farming for traditional medicine
\\(B) Protect Bengal tiger and its habitat through designated tiger reserves
\\(C) Relocate tigers from India to Africa
\\(D) Create a captive breeding program only

\item Match the following protected area categories with their descriptions:\\
1. National Park $\rightarrow$ (a) Strictly protected; no human activity allowed\\
2. Wildlife Sanctuary $\rightarrow$ (b) Some human activities allowed; species-focused protection\\
3. Biosphere Reserve $\rightarrow$ (c) Core, buffer, transition zones; human use in outer zones\\
4. Sacred groves $\rightarrow$ (d) Traditional community-protected forest patches
\\(A) 1-a, 2-b, 3-c, 4-d
\\(B) 1-b, 2-a, 3-d, 4-c
\\(C) 1-c, 2-d, 3-a, 4-b
\\(D) 1-d, 2-c, 3-b, 4-a

\item Genetic diversity is important because it:
\\(A) Determines ecosystem productivity
\\(B) Provides the raw material for evolution and adaptation; reduces vulnerability to disease and environmental change
\\(C) Determines species richness
\\(D) Controls nutrient cycling

\end{enumerate}

\newpage
\title{Chapter 38: Environmental Issues}
\maketitle

\begin{enumerate}[label=Q\arabic*.]

\item The ozone layer in the stratosphere absorbs:
\\(A) Infrared radiation
\\(B) Harmful UV-B and UV-C radiation from the sun
\\(C) Visible light
\\(D) Microwave radiation

\item Chlorofluorocarbons (CFCs) deplete the ozone layer because:
\\(A) They absorb ozone directly
\\(B) UV radiation breaks down CFCs releasing Cl radicals that catalytically destroy O$_3$ (each Cl can destroy 100,000 O$_3$ molecules)
\\(C) They react with N$_2$ in the stratosphere
\\(D) They increase stratospheric temperature

\item The greenhouse effect is caused by greenhouse gases (CO$_2$, CH$_4$, N$_2$O, H$_2$O) because they:
\\(A) Reflect sunlight back to space
\\(B) Absorb and re-emit infrared radiation emitted by Earth's surface, warming the lower atmosphere
\\(C) Block solar radiation from reaching Earth's surface
\\(D) Destroy stratospheric ozone causing warming

\item Acid rain is primarily caused by:
\\(A) CO$_2$ dissolving in rainwater
\\(B) SO$_2$ and NO$_x$ from burning fossil fuels reacting with water to form sulfuric and nitric acids
\\(C) Ozone reacting with water vapor
\\(D) Agricultural fertilizers releasing ammonia

\item Assertion: Eutrophication leads to decrease in dissolved oxygen in water bodies.\\
Reason: Algal blooms from eutrophication block sunlight; when algae die, decomposers consume oxygen, creating hypoxia.
\\(A) Both true, Reason is correct explanation
\\(B) Both true, Reason is not correct explanation
\\(C) Assertion true, Reason false
\\(D) Assertion false, Reason true

\item Biomagnification of DDT results in highest concentration in:
\\(A) Phytoplankton (primary producers)
\\(B) Top predators (e.g., hawks, large fish, humans)
\\(C) Water
\\(D) Primary consumers (herbivores)

\item The Montreal Protocol (1987) was an international agreement to:
\\(A) Reduce greenhouse gas emissions
\\(B) Phase out production and use of ozone-depleting substances (CFCs, halons, methyl bromide)
\\(C) Protect biodiversity
\\(D) Reduce acid rain

\item Chipko movement in India (1970s) was a grassroots environmental movement aimed at:
\\(A) Protesting against water pollution
\\(B) Preventing deforestation; villagers hugged trees to prevent their felling
\\(C) Demanding clean air legislation
\\(D) Promoting renewable energy

\item E-waste (electronic waste) is hazardous because it contains:
\\(A) Only plastic waste
\\(B) Heavy metals and toxic materials (lead, mercury, cadmium, arsenic) that contaminate soil and groundwater
\\(C) Radioactive materials exclusively
\\(D) Organic compounds only

\item The primary source of methane (a potent greenhouse gas) includes:
\\(A) Automobile exhaust
\\(B) Paddy fields, cattle rumen fermentation, landfills, and natural gas leaks
\\(C) Coal burning only
\\(D) Refrigerants

\item Noise pollution above which threshold (dB) is considered harmful to human hearing with prolonged exposure?
\\(A) 30 dB
\\(B) 75 dB (harmful above 80 dB with prolonged exposure)
\\(C) 120 dB (immediate hearing damage)
\\(D) 45 dB

\item Which of the following is NOT a consequence of global warming?
\\(A) Rising sea levels due to glacial melting
\\(B) Increased frequency of extreme weather events
\\(C) Decrease in atmospheric CO$_2$ concentration
\\(D) Shifts in species distribution ranges

\item The Kyoto Protocol (1997) required:
\\(A) Developed nations to reduce biodiversity loss
\\(B) Industrialized nations to reduce greenhouse gas emissions by specific targets
\\(C) All nations to eliminate CFC production
\\(D) Developing nations to reduce deforestation

\item Match the following pollutants with their primary sources:\\
1. SO$_2$ $\rightarrow$ (a) Coal-burning power plants\\
2. CO $\rightarrow$ (b) Incomplete combustion in automobiles\\
3. NO$_x$ $\rightarrow$ (c) High-temperature combustion (vehicles, industry)\\
4. CFCs $\rightarrow$ (d) Refrigerants and aerosol propellants
\\(A) 1-a, 2-b, 3-c, 4-d
\\(B) 1-b, 2-a, 3-d, 4-c
\\(C) 1-c, 2-d, 3-a, 4-b
\\(D) 1-d, 2-c, 3-b, 4-a

\item Integrated Pest Management (IPM) is preferred over chemical pesticides because:
\\(A) IPM uses stronger chemicals
\\(B) IPM combines biological, cultural, physical, and chemical methods to minimize pesticide use, reducing environmental impact
\\(C) IPM is cheaper to implement always
\\(D) IPM eliminates all pests completely

\end{enumerate}
\end{document}
